% ===== §22 Aesthetic Education as the Reduction Terminus =====
\section{The Terminus of the Reduction Chain: The Unification of the Series's Fields}
\label{p22:sec:terminus}

\noindent
The reduction chain that has organised this paper reaches its terminus here, and the terminus has a retrospective force the earlier sections could only anticipate. The chain ran from relational production through the field, from the field's construction as an aesthetic rather than an engineering matter, through the coupling of the aesthetic and the ethical, to the reframing of field-building as aesthetic education. Stated forward, this is the argument of the paper. Stated backward, from the terminus, it does something more: it unifies the whole series, because every paper in the series has described a field, and the capacity this paper names is the capacity to build them all. What looked, across the series, like a sequence of distinct studies, of poetry and the vow and water and travel, is revealed from the terminus as a sequence of studies of fields, and the aesthetic education this paper describes is the cultivation of the single capacity that builds every one of them.

\subsection{Every paper described a field}
\label{p22:sub:terminus-every-field}

The retrospective unification can be made exact by recalling what each of the prior studies was a study of. The study of poetry described a field, the conditional environment in which meaning resists settlement and moves along the interpretive path, a field whose generativity lay in its refusal of closure. The study of the vow described a field, the conditional environment of trust in which the unsecured covenant grounds a coupling more firmly than the secured contract, a field whose generativity lay in its self-legislation. The study of water described a field, the conditional environment of the yielding that overcomes by not contending, a field whose generativity lay in its non-extractive flow. The study of travel described a field, the conditional environment of spaciousness in which contingency is met and the relational subject generated, a field whose generativity lay in its openness. These were not four unrelated topics but four fields, four conditional environments of generative coupling, each studied for the particular way its conditions secured the good cycle, and the variety among them is the variety of conditions under which a field can be generative, not a variety of unrelated subjects.

\claim{Every prior study in the series described a field, a conditional environment of generative coupling, distinguished by the particular conditions under which it secured the good cycle: poetry by its refusal of closure, the vow by its self-legislation, water by its non-extractive flow, travel by its spaciousness. The aesthetic education this paper describes is the cultivation of the single capacity that builds them all, the capacity to perceive a field's state, generate the appropriate act, and sustain the accumulation, in whatever conditions the particular field presents. The terminus of the reduction chain is therefore also the unification of the series: what the series studied as distinct fields, this paper names as the objects of one cultivable capacity.}

The four-step model gives the unification a finer grain, because each of the prior fields can be seen to call on the four movements, weighted differently as its conditions require. The field of poetry calls especially on the freshness of reproduction, the refusal of closure being the refusal to let the cycle settle into a fixed reading, the keeping-open that is reproduction guarding against its degenerate form; to build the field of poetry is to reproduce without ever closing, the fourth movement held in the mode of perpetual non-settlement. The field of the vow calls especially on reception, the self-legislation being the receiving of a standard one then holds oneself to, the taking-in of an obligation that binds because it has been let in; to build the field of the vow is to receive in the strong sense, to let in what will bind. The field of water calls especially on the non-grasping that runs through all four, the yielding that overcomes by not contending being the non-extractive flow that the whole model requires, reception that does not seize and reproduction that does not hoard; to build the field of water is to perform the four movements in the mode of yielding. The field of travel calls especially on perception, the spaciousness being the openness in which contingency is perceived, the room left for the unplanned to be seen; to build the field of travel is to perceive in the mode of openness to contingency. Each prior field is thus the four-step model weighted toward one of its movements by the field's particular conditions, and the single capacity builds them all by performing the four movements with the weighting each field's conditions require, which is the finer form of the unification, the one model differently weighted across the series' several fields.

\subsection{One capacity, many conditions}
\label{p22:sub:terminus-one-capacity}

The unification is not a flattening, and the section must be careful to preserve the variety it gathers. The fields of the series differ really, in the conditions under which each secures its generativity, and the capacity to build them is not a single technique applied uniformly but a single power exercised differently as the conditions differ. To build the field of poetry is to sustain the refusal of closure; to build the field of the vow is to perform the self-legislation; to build the field of water is to yield without contending; to build the field of travel is to keep the spaciousness open. These are different exercises of one capacity, the capacity to perceive what a field's generativity requires in its particular conditions and to act to secure it, and the unity is in the capacity and not in a uniform method. This is why the aesthetic education this paper describes is an education of a perception and not a transmission of a technique: a technique would apply uniformly and would fail across the variety of fields, whereas a trained perception discerns, in each field, what its particular conditions require, and acts accordingly. The single capacity builds many fields not by doing the same thing in each but by perceiving, in each, what that field's generativity demands.

This is also why the present paper completes rather than merely extends the series. The prior studies each described the conditions of a particular field and the way those conditions secured the good cycle; none described the capacity to build such fields as such, the perception-and-practice that, trained, can secure generativity across the variety of conditions. This paper describes that capacity, and in describing it gives the series its missing member, the account of the cultivation that the prior studies presupposed without naming. When the study of poetry showed how a poem sustains the good cycle, it presupposed a capacity to sustain it; when the study of the vow showed how a covenant grounds trust, it presupposed a capacity to perform the grounding; and so through the series, each study presupposing the cultivable capacity that this paper finally names. The aesthetic education of the everyday is thus the series' keystone in the architectural sense, the member that the others leaned toward and that, set in place, holds them together, the account of how a person comes to be able to build the generative fields that the whole series has studied.

\subsection{The everyday, the privileged site}
\label{p22:sub:terminus-everyday}

It remains to say why this capstone capacity is cultivated in the everyday rather than in the exceptional fields the series often studied, the heightened field of the poem or the solemn field of the vow. The answer is that the everyday is where the capacity is most generally needed and most generally available, and a cultivation of the capacity that required the exceptional field would be a cultivation available only to those who could enter exceptional fields. The poem and the vow and the journey are not daily; they are the heightened occasions on which the good cycle is most visible, and the series studied them partly because their heightening made the structure of the good cycle legible. But a life is mostly not heightened, mostly daily, mostly the unspacious repetition that no journey relieves, and a capacity to build generative fields that worked only on the heightened occasions would leave the bulk of a life uncultivated. The everyday is the privileged site of aesthetic education because it is where the capacity must work if it is to work at all for most of a life, and because, as the paper has argued throughout, the capacity can work there, the good cycle being buildable in the narrow room of the daily and not only in the spacious occasions of the exceptional. The series studied the heightened fields to make the good cycle legible; this paper returns the capacity to the daily field, where it is needed most and where, the paper has shown, it is genuinely available, completing the series by bringing its central capacity home to the everyday from which the exceptional fields were always only the heightened departures.

\subsection{The choice of the name aesthetic education}
\label{p22:sub:terminus-name}

A final objection concerns the name. If the capacity is the capacity to build generative relational fields, and if the building is at once aesthetic and ethical, why call its cultivation aesthetic education rather than ethical education, or relational education, or by some new name that would not import the museum's associations the introduction worked to leave behind? The choice of name is deliberate and the section defends it. The cultivation is not called ethical education because that name would suggest the transmission of a moral code, the very rule-governed instruction the paper has argued the appropriate exceeds; ethical education, as ordinarily understood, teaches principles, and the capacity this paper describes is precisely not the application of principles but the trained perception of the fitting where no principle decides. It is not called relational education because that name would lose the central thesis, that the building of the field is a matter of the perception of appropriateness, which is an aesthetic perception in the broad sense the genealogy recovered, the discernment of the fitting where no rule applies. The name aesthetic education is chosen because it alone names the capacity correctly, as the cultivation of a perception of the fitting, while the genealogy and the coupling have shown that this aesthetic perception is at the same time the perception of the good, so that aesthetic education, rightly understood, already is ethical education, not by adding a moral code but by cultivating the perception in which the appropriate and the good are one. The name is not a metaphor or a borrowing; it is the precise name for a cultivation of perception that the tradition narrowed to the museum and the paper has returned to its full breadth, the breadth in which to educate the aesthetic sensibility is to educate the moral one, because they were never two. To call it aesthetic education is to insist on exactly the coupling the paper established, and to call it anything else would be to lose the coupling in the naming.
